% Five-stage pipeline (sample figure of Lesson 1).
% Shows the IF, ID, EX, MEM and WB stages as boxes joined by arrows, each
% with its name below; the names switch with the document language.
% Usage: % Five-stage pipeline (sample figure of Lesson 1).
% Shows the IF, ID, EX, MEM and WB stages as boxes joined by arrows, each
% with its name below; the names switch with the document language.
% Usage: % Five-stage pipeline (sample figure of Lesson 1).
% Shows the IF, ID, EX, MEM and WB stages as boxes joined by arrows, each
% with its name below; the names switch with the document language.
% Usage: % Five-stage pipeline (sample figure of Lesson 1).
% Shows the IF, ID, EX, MEM and WB stages as boxes joined by arrows, each
% with its name below; the names switch with the document language.
% Usage: \input{figures/l01-pipeline-stages}
% Width: body (120 mm); the drawing is about 96 mm wide.
\begin{tikzpicture}[
  stage/.style={draw, rounded corners=1pt, minimum width=1.6cm, minimum height=0.8cm,
                font=\sffamily\small, fill=ln-box},
  >={Stealth[length=2mm]}, node distance=4mm,
]
  \node[stage] (if)  {IF};
  \node[stage, right=of if]  (id)  {ID};
  \node[stage, right=of id]  (ex)  {EX};
  \node[stage, right=of ex]  (mem) {MEM};
  \node[stage, right=of mem] (wb)  {WB};
  \draw[->] (if) -- (id);
  \draw[->] (id) -- (ex);
  \draw[->] (ex) -- (mem);
  \draw[->] (mem) -- (wb);
  \node[below=1mm of if,  font=\footnotesize\color{ln-muted}] {\iflangja{命令フェッチ}{fetch}};
  \node[below=1mm of id,  font=\footnotesize\color{ln-muted}] {\iflangja{デコード}{decode}};
  \node[below=1mm of ex,  font=\footnotesize\color{ln-muted}] {\iflangja{実行}{execute}};
  \node[below=1mm of mem, font=\footnotesize\color{ln-muted}] {\iflangja{メモリ}{memory}};
  \node[below=1mm of wb,  font=\footnotesize\color{ln-muted}] {\iflangja{書き戻し}{write back}};
\end{tikzpicture}

% Width: body (120 mm); the drawing is about 96 mm wide.
\begin{tikzpicture}[
  stage/.style={draw, rounded corners=1pt, minimum width=1.6cm, minimum height=0.8cm,
                font=\sffamily\small, fill=ln-box},
  >={Stealth[length=2mm]}, node distance=4mm,
]
  \node[stage] (if)  {IF};
  \node[stage, right=of if]  (id)  {ID};
  \node[stage, right=of id]  (ex)  {EX};
  \node[stage, right=of ex]  (mem) {MEM};
  \node[stage, right=of mem] (wb)  {WB};
  \draw[->] (if) -- (id);
  \draw[->] (id) -- (ex);
  \draw[->] (ex) -- (mem);
  \draw[->] (mem) -- (wb);
  \node[below=1mm of if,  font=\footnotesize\color{ln-muted}] {\iflangja{命令フェッチ}{fetch}};
  \node[below=1mm of id,  font=\footnotesize\color{ln-muted}] {\iflangja{デコード}{decode}};
  \node[below=1mm of ex,  font=\footnotesize\color{ln-muted}] {\iflangja{実行}{execute}};
  \node[below=1mm of mem, font=\footnotesize\color{ln-muted}] {\iflangja{メモリ}{memory}};
  \node[below=1mm of wb,  font=\footnotesize\color{ln-muted}] {\iflangja{書き戻し}{write back}};
\end{tikzpicture}

% Width: body (120 mm); the drawing is about 96 mm wide.
\begin{tikzpicture}[
  stage/.style={draw, rounded corners=1pt, minimum width=1.6cm, minimum height=0.8cm,
                font=\sffamily\small, fill=ln-box},
  >={Stealth[length=2mm]}, node distance=4mm,
]
  \node[stage] (if)  {IF};
  \node[stage, right=of if]  (id)  {ID};
  \node[stage, right=of id]  (ex)  {EX};
  \node[stage, right=of ex]  (mem) {MEM};
  \node[stage, right=of mem] (wb)  {WB};
  \draw[->] (if) -- (id);
  \draw[->] (id) -- (ex);
  \draw[->] (ex) -- (mem);
  \draw[->] (mem) -- (wb);
  \node[below=1mm of if,  font=\footnotesize\color{ln-muted}] {\iflangja{命令フェッチ}{fetch}};
  \node[below=1mm of id,  font=\footnotesize\color{ln-muted}] {\iflangja{デコード}{decode}};
  \node[below=1mm of ex,  font=\footnotesize\color{ln-muted}] {\iflangja{実行}{execute}};
  \node[below=1mm of mem, font=\footnotesize\color{ln-muted}] {\iflangja{メモリ}{memory}};
  \node[below=1mm of wb,  font=\footnotesize\color{ln-muted}] {\iflangja{書き戻し}{write back}};
\end{tikzpicture}

% Width: body (120 mm); the drawing is about 96 mm wide.
\begin{tikzpicture}[
  stage/.style={draw, rounded corners=1pt, minimum width=1.6cm, minimum height=0.8cm,
                font=\sffamily\small, fill=ln-box},
  >={Stealth[length=2mm]}, node distance=4mm,
]
  \node[stage] (if)  {IF};
  \node[stage, right=of if]  (id)  {ID};
  \node[stage, right=of id]  (ex)  {EX};
  \node[stage, right=of ex]  (mem) {MEM};
  \node[stage, right=of mem] (wb)  {WB};
  \draw[->] (if) -- (id);
  \draw[->] (id) -- (ex);
  \draw[->] (ex) -- (mem);
  \draw[->] (mem) -- (wb);
  \node[below=1mm of if,  font=\footnotesize\color{ln-muted}] {\iflangja{命令フェッチ}{fetch}};
  \node[below=1mm of id,  font=\footnotesize\color{ln-muted}] {\iflangja{デコード}{decode}};
  \node[below=1mm of ex,  font=\footnotesize\color{ln-muted}] {\iflangja{実行}{execute}};
  \node[below=1mm of mem, font=\footnotesize\color{ln-muted}] {\iflangja{メモリ}{memory}};
  \node[below=1mm of wb,  font=\footnotesize\color{ln-muted}] {\iflangja{書き戻し}{write back}};
\end{tikzpicture}
